\subsection*{NS-50: exact fixed-width localization and the critical remainder}

\paragraph{Status and provenance.}
This is an analytic attempt at the complete uniform signed floor, not a
proof of CAE, G2 or RH. It uses the v1.51 complete form, the already
certified positive window $\lambda=4$, and
Lemma~\ref{lem:v136-total-radicals}. The localization identity below is
derived directly. The failed remainder estimate is a specialization of
NS-47's ordinary-closable-comparison obstruction; its proof is included
so this checkpoint needs no pending NS-47 manuscript integration. No
new window, tail metric, numerical illustration or interval gate is used.

\paragraph{The signed form and the localization.}
Put $\mathcal C=C_c^\infty(\mathbb R)$ and use zero extension and the
unitary Fourier transform throughout. For $f\in\mathcal C$ write
\[
 c_f(t)=\int_{\mathbb R} f(x+t)\overline{f(x)}\,dx,
 \qquad w_n=\frac{\Lambda(n)}{\sqrt n},\qquad \ell_n=\log n,
 \qquad \rho(t)=\frac{e^{t/2}}{2\sinh t}.
\]
The complete support-consistent form in
Proposition~\ref{prop:v135-both-parities} is equivalently
\begin{align}
 q[f]={}&c_{\rm ar}\norm f^2
       +2\int_0^\infty\rho(t)
                   (\norm f^2-\Re c_f(t))\,dt\notag\\
 &-2\sum_{n\ge2}w_n\Re c_f(\ell_n)
       +4\int_0^\infty\cosh(t/2)\Re c_f(t)\,dt,
 \label{eq:ns50-q-correlation}
\end{align}
where $c_{\rm ar}=\psi(1/4)-\log\pi$. For each compact test the
prime sum is finite, and its strict endpoint convention is automatic
because $c_f(t)=0$ at and beyond the support diameter. The pole is the
full kernel $2\cosh((x-y)/2)$; its negative odd square has not been
discarded. The archimedean multiplier is
$\alpha(\xi)=\Re\psi(1/4+i\xi/2)-\log\pi$.

Fix a nonzero real even $\chi\in C_c^\infty(-b,b)$ with
$0<b<\log4$ and $\norm\chi_2=1$. Set
\[
 \chi_s(x)=\chi(x-s),\qquad
 H(t)=\int_{\mathbb R}\chi(x+t)\chi(x)\,dx,
 \qquad V(t)=1-H(t),
\]
and define
\[
 \mathcal L_\chi[f]=\int_{\mathbb R}q[\chi_s f],ds,
 \qquad \mathcal D_\chi[f]=\mathcal L_\chi[f]-q[f].
\]
Here $\mathcal D_\chi$ denotes a signed remainder, not a positive
Dirichlet energy. We have $H(0)=1$, $H$ real and even,
$\operatorname{supp}H\subset[-2b,2b]$, $V\ge0$, and
\[
 V(t)=\tfrac12\norm{\chi(\cdot+t)-\chi}^2
       \le\tfrac12t^2\norm{\chi'}^2.
\]
Translations preserve \eqref{eq:ns50-q-correlation}. After recentering,
each $\chi_s f$ lies inside the already positive $\lambda=4$ window.
Thus $\mathcal L_\chi[f]\ge0$. This uses both certified parities;
a translated localized piece need not retain the parity or the source
orthogonality of $f$.

\begin{proposition}[Exact signed localization remainder]
\label{prop:ns50-localization}
For every $f\in\mathcal C$,
\begin{align}
 \mathcal D_\chi[f]={}&
  2\int_0^\infty \rho(t)V(t)\Re c_f(t)\,dt
  +2\sum_{n\ge2}w_n V(\ell_n)\Re c_f(\ell_n)\notag\\
 &-4\int_0^\infty\cosh(t/2)V(t)\Re c_f(t)\,dt.
 \label{eq:ns50-remainder}
\end{align}
In particular all signed prime, archimedean and pole cross terms are
retained. At every fixed complete window $I_a$, this remainder extends
to a bounded Hermitian form in ordinary $L^2(I_a)$, with a bound that
may depend on $a$.

Moreover
\begin{equation}\label{eq:ns50-multiplier}
 \mathcal L_\chi[f]=\int_{\mathbb R}
                m_\chi(\xi)|\widehat f(\xi)|^2\,d\xi,
 \qquad m_\chi(\xi)=q[e^{i\xi x}\chi(x)]\ge0,
\end{equation}
and the continuous real even multiplier satisfies
\begin{equation}\label{eq:ns50-high-frequency}
 m_\chi(\xi)=\log(2+|\xi|)+O_\chi(1).
\end{equation}
Consequently $\mathcal L_\chi$ has a closed nonnegative extension
on the global logarithmic Fourier form domain; its square-root
multiplier $T_\chi f=\sqrt{m_\chi}\widehat f$ is closed and unbounded.
These statements also hold in either parity sector.
\end{proposition}
\begin{proof}
Fubini and the normalization of $\chi$ give
\[
 \int\norm{\chi_s f}^2\,ds=\norm f^2,
 \qquad \int c_{\chi_s f}(t)\,ds=H(t)c_f(t).
\]
Substitute these two identities into
\eqref{eq:ns50-q-correlation}. The difference of the two
archimedean terms is $2\int\rho V\Re c_f$, the difference of the
prime terms is $2\sum w_nV(\ell_n)\Re c_f(\ell_n)$, and the pole
difference has the opposite sign. This proves
\eqref{eq:ns50-remainder}. Near zero, $\rho(t)=O(1/t)$ and
$V(t)=O_\chi(t^2)$; at infinity $\rho$ decays exponentially.
Thus $\rho V\in L^1(0,\infty)$. The prime and pole correlations
have compact $t$ support for a compact $f$, so every interchange is
absolutely justified after the archimedean difference is grouped.
The original archimedean energy can also be integrated as a
nonnegative squared difference before this subtraction.

For $\operatorname{supp}f\subset I_a$, Cauchy--Schwarz gives the
explicit fixed-window bound
\begin{align*}
 |\mathcal D_\chi[f]|\le\Bigl(&2\int_0^\infty\rho V
       +2\sum_{n<e^{2a}}w_nV(\ell_n)
       +4\int_0^{2a}\cosh(t/2)V(t)\,dt\Bigr)\norm f^2.
\end{align*}
This is only a convergence/domain bound. Its unsigned terms are not
proposed as a cofinal estimate.

Taking Fourier transforms of the localized correlation yields
\begin{align*}
 m_\chi(\xi)={}&\alpha(\xi)
       +2\int_0^\infty\rho(t)V(t)\cos(\xi t)\,dt\notag\\
 &-2\sum_{n\ge2}w_nH(\ell_n)\cos(\xi\ell_n)
       +4\int_0^\infty\cosh(t/2)H(t)\cos(\xi t)\,dt.
\end{align*}
The last three terms are bounded uniformly in $\xi$: the first
integrand is integrable, the prime sum has finitely many terms, and
the last integrand has compact support. The digamma asymptotic
proves \eqref{eq:ns50-high-frequency}. The same formula equals
$q[e^{i\xi x}\chi]$, since its correlation is $e^{i\xi t}H(t)$.
Fixed-window positivity proves its nonnegativity for every real
$\xi$. Continuity and evenness follow either from this expression
or from the multiplier formula. A nonnegative measurable multiplier
defines a closed square-root multiplication operator. The two-sided
logarithmic estimate identifies its form domain and proves
unboundedness. Evenness of $m_\chi$ makes it commute with reflection;
both its even and odd restrictions are unbounded, for example by
normalized cosine and sine modulations of a fixed even smooth bump.
On each fixed window, closure from the physical smooth core extends
the identities to the complete logarithmic form domain.
\end{proof}

\begin{proposition}[Bounded or strictly absorbable localization errors fail]
\label{prop:ns50-critical-coefficient}
Fix $0\le\theta<1$. There is no finite $C$ such that
\begin{equation}\label{eq:ns50-failed-bound}
 \mathcal D_\chi[f]\le
            \theta\mathcal L_\chi[f]+C\norm f^2
 \quad\hbox{for all }f\in\mathcal C.
\end{equation}
The same failure holds separately on the full even core and on the
odd core. In particular the signed localization error has no
window-independent upper ordinary-norm bound, including under RH.
This excludes this stronger sufficient comparison, not CAE.

For either full parity core let
\[
 B_{\theta,\chi}^{\pm}(a)=
 \sup_{\substack{f\in C_c^\infty(I_a),\ f\ \text{of parity}\ \pm\\
                 \norm f=1}}
   \bigl(\mathcal D_\chi[f]-\theta\mathcal L_\chi[f]\bigr).
\]
These finite fixed-window quantities increase to $+\infty$ as
$a\to\infty$. No rate or sign of $q_a$ is inferred.
\end{proposition}
\begin{proof}
The proposed estimate is exactly
\[
 q[f]\ge(1-\theta)\mathcal L_\chi[f]-C\norm f^2.
\]
Let $\mathcal R$ be the dense linear translate span in
Lemma~\ref{lem:v136-total-radicals}. For every fixed $r\in\mathcal R$
there are compact cutoffs $r_a\to r$ in ordinary norm with
$q[r_a]\to0$, by the complete ordinary operator residual.
With $T=\sqrt{1-\theta}\,T_\chi$, the estimate gives
\[
 \limsup_a\norm{Tr_a}^2\le C\norm r^2.
\]
Choose a weakly convergent subsequence of the bounded images. Since
the graph of the closed linear operator $T$ is a norm-closed linear
subspace, it is weakly closed. The pairs $(r_a,Tr_a)$ therefore have
a weak graph limit $(r,v)$, and
$r\in\mathcal D(T)$ with $\norm{Tr}^2\le C\norm r^2$.
This holds on the dense linear space $\mathcal R$, including all
differences. Approximating any vector by that space and using
closedness again makes $T$ everywhere defined with norm at most
$\sqrt C$. This contradicts \eqref{eq:ns50-high-frequency}.
For either parity use the reflected dense translate span and its
same-parity cutoffs from the same lemma.

The preceding proposition bounds $\mathcal D_\chi$ on each fixed
window, and $\theta\mathcal L_\chi\ge0$, so each displayed supremum
is finite. Support consistency makes it nondecreasing. A finite
cofinal upper bound would imply \eqref{eq:ns50-failed-bound} on
the entire corresponding compact parity core, which has just been
excluded. It therefore diverges. No assumption on RH has entered.
\end{proof}

\paragraph{What the source condition does and does not change.}
The odd core is already source admissible. Its obstruction alone
prevents using \eqref{eq:ns50-failed-bound} with a common bounded
error in both required sectors. No density or monotonicity assertion
is made here for the varying even spaces $u_a^\perp$; the full-even
statement must not be silently renamed a source-complement statement.
Localization need not preserve that constraint on each piece because
the fixed-window positivity input is for the full two-parity form.
The eventual target is still tested on the original admissible $f$.

\paragraph{The critical signed input remains open.}
For the coefficient exactly one, the inequality
\begin{equation}\label{eq:ns50-open-critical}
 \mathcal D_\chi[f]\le\mathcal L_\chi[f]+C\norm f^2
\end{equation}
is exactly $q[f]\ge-C\norm f^2$. On every even
$f\perp u_a$ and every odd $f$ in a cofinal family of complete
domains, with the same finite $C$, it is therefore exactly CAE.
The established source residual and
Corollary~\ref{cor:v136-complement-floor} would transfer it to a
full floor and RH. Equation~\eqref{eq:ns50-open-critical} is not
proved here. It retains the signed prime-pole cancellation of
\eqref{eq:ns50-remainder}; bounding each component separately would
change the task.

A bounded Fourier cap $m_{\chi,M}=\min(m_\chi,M)$ removes the
unbounded-factor obstruction but supplies no estimate: putting
$\mathcal L_{\chi,M}=\int m_{\chi,M}|\widehat f|^2$ and
$\mathcal D_{\chi,M}=\mathcal L_{\chi,M}-q$ again leaves precisely
$\mathcal D_{\chi,M}\le\mathcal L_{\chi,M}+C\norm f^2$.
For a fixed $M$, the stronger-looking bounded remainder condition
$\mathcal D_{\chi,M}\le C_M\norm f^2$ is equivalent to existence
of a full uniform floor, up to changing constants: a floor $-C$
gives $C_M=M+C$, while such a $C_M$ gives the floor $-C_M$.
Thus frequency capping avoids the proved exclusion but proves no
new arithmetic statement. Likewise taking a coefficient greater
than one in \eqref{eq:ns50-failed-bound} leaves a negative multiple
of the unbounded $\mathcal L_\chi$, hence does not by itself yield
the desired ordinary floor.

The attempted construction had to change from bounded localization
error to critical signed absorption. The former fails for the
specified fixed-width averaging, even though every local piece is
positive. The latter is a new expression of the old missing estimate,
not progress on its truth. Window-dependent localization, other
completions and other arithmetic constructions are not excluded by
this result. No failure of the physical Weil form is asserted.
